\section{Firmware Update}\label{firmware-update}

Firmware is updated from the SD card by the bootloader.

\draftnote{The project documents describe the update at build level only. Confirm the file name,
the bootloader messages and the recovery behaviour with the project owner and replace the text
below with what the device actually shows.}

\subsection{Update procedure}\label{update-procedure}

\begin{enumerate}
\def\labelenumi{\arabic{enumi}.}
\tightlist
\item
  Make sure the SD card uses \textbf{FAT32}.
\item
  Copy the firmware image, \texttt{LXRV2\_lxr02.img}, to the
  \textbf{root folder} of the card.
  \verify{Image file name in the released package.}
\item
  Insert the card into the powered-off LXR.
\item
  Push and hold the \textbf{main encoder}, then switch the power on.
\end{enumerate}

The device starts the bootloader instead of the application. The
bootloader reads the image from the card, checks the image header and
checksum, and writes the application. When it has finished, power the
device off and on again.

\oledsingle{Bootloader}

\verify{Bootloader display texts for ``updating'', ``success'' and the error cases. The previous product showed \texttt{Bootloader v1.1}, \texttt{updating... (1/2)} and \texttt{success! please reboot...}.}

\figplaceholder{firmware_update_steps.png}{Entering the bootloader.}{Three panels: SD card inserted in the slot, a finger holding the main encoder, power switch turned on. Add the expected display text below the third panel.}

\subsection{Image format}\label{image-format}

The image is a binary file with an 8-byte marker \texttt{LXRV2IMG}, the
payload size, a checksum and then the program. The bootloader occupies
the first 32 KB of flash and the application starts above it. The area
at the top of flash (five 256 KB sectors) is reserved for samples and is
never touched by a firmware update.

\subsection{After an update}\label{after-an-update}

\begin{itemize}
\tightlist
\item
  A flash image check verifies a stored checksum for every program
  sector. If a sector is wrong the display shows
  \texttt{Img\ BAD\ s:...} and waits for BAR1.
\item
  A device that was updated from an older firmware may report the sample
  catalogue as stale; reinstall the samples (see the Storage chapter).
\end{itemize}

\subsection{Rolling back}\label{rolling-back}

\attention{Firmware older than the current sample layout plays program code as audio if it reads an
enlarged image. Reinstall samples at low volume before rolling back to a version older than the
one that introduced the current layout.}

\verify{The exact firmware version boundary for the rollback hazard (the engineering notes name Session 073).}

\subsection{Bootloader error messages}\label{bootloader-error-messages}

\verify{Replace with the real list. The previous product reported: Firmware error (image not found or card not FAT32), Header error (corrupt image), EOF error (truncated image) and Mainboard error.}
